\chapter*{Executive summary}
\addcontentsline{toc}{chapter}{Executive summary}
\markboth{Executive summary}{}

Under GPS jamming, a soldier often doesn't know where he is. He can still say what he sees:
``the castle keep at my 11 o'clock, a church at bearing 240, sun at my 1 o'clock''. Near turns
that sentence into a position. When the position is still vague, it asks the one spoken question
that narrows it down most. Then it reads the answer back, and puts the casualty on the
coordinator's map with a MIST card, a tourniquet timer and a place in the evacuation order.
The name of the method is old: the \emph{problème de Pothenot} is the French surveyors' name for
finding your own position from the angles to known points.

Around eleven teams took the same hackathon challenge. All of them map what a speaker says, and
all of them assume the speaker knows where he is. We built for the case where he doesn't.

\paragraph{What Near does, end to end.}
\begin{itemize}
  \item \textbf{Locates the speaker without GPS.} A likelihood grid (25 m cells $\times$ 72 headings)
        combines spoken clock directions, compass bearings, the sun, landmarks seen or not seen
        (with a LiDAR line-of-sight model), ground clues, sounds, echoes and the noise level.
  \item \textbf{Asks back.} Candidate questions are scored by the position error and the entropy they
        are expected to remove. Fewer and easier questions are asked when the voice sounds stressed.
  \item \textbf{Organises the casualty.} MIST card, MARCH state, ``worse in'' countdowns, mobility,
        an impact graph from wound to kit, team size, a 3D body twin and an optional heart rate
        from a chest recording. All of this is decision support only.
  \item \textbf{Maps threats.} Drones, gunfire and explosions are heard by the microphone, or named
        by voice, and drawn as zones; overlapping reports from several units mark the likely
        source area.
  \item \textbf{Brings people back.} A casualty who can walk is guided to the collection point along
        the safest walking route, with spoken landmark-based legs and a pin that moves by dead
        reckoning. A casualty who can't walk gets a medic route instead.
  \item \textbf{Works in war conditions.} With no internet, no GPS and no phone network, everything
        runs on one laptop: local Whisper speech-to-text, a local Mistral 7B, offline map tiles,
        offline voices. The laptop listens to the tactical voice radio, routes each message by
        callsign and answers on the radio. A phone, if the unit has one, runs a 5.5M-parameter
        sound model offline and sends its measurements as ~2 s of audio tones through the same
        radio.
\end{itemize}

\paragraph{Key numbers (all measured; methods in Chapter~\ref{ch:eval}).}
\begin{itemize}
  \item Spots where at least two landmarks are truly visible, ask-back loop with up to three
        questions: median error \textbf{43 m}.
  \item Spoken compass bearings to two named landmarks: median error \textbf{71 m}, against
        250 m with clock directions.
  \item Aim \& tap (phone tilt and turn): median error \textbf{75 m} in the streets where a landmark
        top is in sight, against 417 m before the 3D model.
  \item Drones heard through a simulated radio: \textbf{100 \%} detected by our radio-trained
        5.5M-parameter model, against 6--9 \% for the 86M-parameter drone detector it learned from.
  \item Data packets through a simulated lossy digital radio: \textbf{39/40} delivered with
        two-copy combining, against 10/40 with one copy.
  \item Safest route on the city's walking network: about 30 ms. In the test scene, a 390 m detour
        avoided 398 m inside a small-arms zone.
\end{itemize}

\paragraph{Business.} Near is software. It runs on a laptop and radios the units already own, and
it needs no new hardware in the field. The first customers are military medical services and
command posts in Europe and Ukraine, where GPS interference is routine and drones have made the
``golden hour'' unreachable. The same engine fits civilian rescue: lost callers who can't give
an address, mountain rescue, and disaster zones without a network. Chapter~\ref{ch:business}
covers the market, the competition, the business model and the risks.

\paragraph{What is not proven yet.} Every result above comes from real map data and real recordings, but the soldier,
the radio channel or both are simulated. Before anyone relies on Near, it needs one walk by the
Château de Vincennes with GPS off, recordings through the unit's real radio, and medical
validation of every threshold. Chapter~\ref{ch:roadmap} lists these tests.
